\section{Related Works}
\subsection{LLM Reasoning with RL}
Reinforcement learning has become an essential technique for incentivizing chain-of-thought (CoT) reasoning in LLMs \cite{RL4LRM, Survey-LLMs}.
A milestone paradigm is reinforcement learning with verifiable rewards (RLVR) \cite{DeepSeekMath, O1, Deepseek-R1}, where LLMs are trained on verifiable tasks such as math problems \cite{AIME}, puzzle solving \cite{ARC-AGI-2}, and code execution \cite{LiveCodeBench}, by rewarding final-answer correctness using task-specific verifiers.
Since such tasks are intrinsically difficult to reward-hack, RLVR tends to encourage LLMs to produce longer CoTs before answering, thereby improving reasoning performance \cite{Deepseek-R1, O1}.
However, despite its effectiveness in incentivizing reasoning, most RLVR training remains outcome-centric and provides no direct supervision over the CoT itself, lacking guidance on how the model should reason.
While several efforts have tried to mitigate this by assigning process rewards using additional reward models (RM) \cite{PRM8K}, training such RMs requires labor-intensive training data collection \cite{PRM8K, PRIME, ProcessBench}, and a static RM often fails to keep pace with the evolving CoT distribution during training \cite{Deepseek-R1}.
Therefore, such challenges hinder the combination of RLVR training with CoT supervision criteria for rewarding in practice.

\subsection{Self-Evolving in LLMs}
A critical step toward general artificial intelligence lies in the ability to self-evolve with little or no human intervention \cite{self-evolving-survey}. 
Recent work has made progress via self-training and self-play, where a single policy model $\pi_\theta$ can self-generate training data \cite{AbsoluteZero, SSP, SPICE}, or self-provide reward signals during learning \cite{SPIN, SeRL, SRLM}. 
The self-evolving capabilities can further be strengthened with techniques such as multi-role reinforcement learning, where a single policy model $\pi_\theta$ is prompted to play different roles (\eg reasoner, data generator, or reward provider) and trained with role-specific rewards \cite{AbsoluteZero, SPICE, Alignmen-Waltz}.
Therefore, these advances in such self-evolving methods suggest the possibilities of self-proposing criteria for CoT supervision, while this direction remains largely under-explored.

\subsection{Reinforcement Learning with Rubrics}
Recently, reinforcement learning with rubrics has emerged as a promising way to provide fine-grained and interpretable reward signals \cite{Rubicon, Rubrics-as-Rewards, HealthBench, Chasing-the-Tail}.
Here, rubrics are structured, checkable evaluation criteria (\ie explicit requirements for assessing outputs) \cite{Rubrics-as-Rewards} that decompose high-level objectives into multi-facet supervision signals \cite{DR-Tulu}, enabling RL to extend towards non-verifiable and open-ended domains \cite{RuscaRL, Rubrics-as-Rewards}.
However, existing work largely applies rubrics to score response quality, leaving their effectiveness for CoT supervision underexplored.
Moreover, most approaches rely on predefined or static rubrics, which may become mismatched as model behaviors evolve during training \cite{DR-Tulu}.
To address these gaps, we introduce a framework that self-proposes rubrics as CoT supervision criteria and continually evolves them throughout training, integrating rubric-based RL for CoT supervision with self-evolving training.


\section{Preliminaries}
In this section, we briefly review the necessary background that forms the foundation of our work, including incentivizing LLM reasoning with RL in Section \ref{sec:pre-training} and multi-role RL under a single policy in Section \ref{sec:pre-masrl}. 

    \subsection{Incentivizing LLM Reasoning with RL} \label{sec:pre-training}
In this work, we focus on incentivizing LLMs to reason with RL, where the policy model is optimized by rewarding the high-quality rollouts.

\textbf{Rollout.}
Given a question $\mathcal{Q}$, the policy model $\pi_{\theta}$ attempts to solve it by generating a response that contains a CoT $\hat{\mathcal{C}}$ followed by an answer $\hat{\mathcal{A}}$, which can be formulated as:
\begin{equation}
\hat{\mathcal{C}}, \hat{\mathcal{A}} \sim \pi_{\theta}(\cdot \mid\mathcal{Q}).
\end{equation}
\textbf{Rewarding.}
After rollouts generation, rewards are assigned based on response quality, which can depend on: (i) the CoT $\hat{\mathcal{C}}$ and (ii) the answer $\hat{\mathcal{A}}$.
In the widely adopted RLVR paradigm, rewards typically depend solely on the correctness of the predicted answer $\hat{\mathcal{A}}$:
\begin{equation}\label{eq:outcome_reward}
r = \psi( \mathbb{I}(\mathcal{A},\hat{\mathcal{A}})),
\end{equation}
where $\mathbb{I}(\cdot,\cdot)$ is an answer correctness indication function which returns $1$ if two answers $\mathcal{A}$ and $\hat{\mathcal{A}}$ are equivalent and $0$ otherwise, and a reward assignment function $\psi$ maps this correctness signal $\mathbb{I}(\mathcal{A},\hat{\mathcal{A}})$ to a binary reward $r$.
Through such RLVR training, the policy model $\pi_\theta$ gradually learns to answer questions, which is often associated with producing longer CoTs \cite{Deepseek-R1}.
% thereby achieving higher answer accuracy.
However, under the RLVR paradigm, the CoT $\hat{\mathcal{C}}$ itself is rarely rewarded directly due to practical implementation challenges \cite{Deepseek-R1}, while the lack of such supervision may induce sub-optimal reasoning strategies \cite{PRM8K}.

\subsection{Multi-Role RL Under a Single Policy Model} \label{sec:pre-masrl}
Recently, studies have shown that a single policy model $\pi_\theta$ can effectively learn multiple roles by assigning role-specific rewards \cite{AbsoluteZero, SPIRAL}:
\begin{equation}\label{eq:reward-role}
r^{\text{role}_i} \rightarrow \pi^{\text{role}_i}_\theta,
\end{equation}
where $r^{\text{role}_i}$ denotes the reward assigned to $\text{role}_i$, and $\pi^{\text{role}_i}_\theta$ represents the same underlying policy model $\pi_\theta$ instantiated under the prompt $\mathcal{P}^{\text{role}_i}$ corresponding to that role.

A typical application under this framework is to let the model simultaneously act as a reasoner and a data generator, achieving RL training starting from no human-labeled data \cite{AbsoluteZero, AlphaZero, SSP}. 
More broadly, this multi-role, single-policy setup provides a natural substrate for self-evolving training: different roles instantiated from the same model can compete \cite{AbsoluteZero, SPICE} or collaborate \cite{Alignmen-Waltz, SPIRAL}, continuously improving the learning signals and behaviors \cite{MAS-Survey}.
Building on this view, we investigate whether the same policy can simultaneously serve as a reasoner and a CoT rubrics generator (\ie rubricator). 
With role-specific rewards, the rubricator is encouraged to produce useful CoT supervision criteria, which in turn enables the self-evolution of CoT rewarding. 
In this case, the reasoner and rubricator collaborate together for better reasoning.


